\section{Proposal II: a branching norm bound from one vector}
\label{sec:branch}

There is relevant prior code here too. The compact edition's
\compact{Lib/GLn/Branching/Branching}{\code{exists_branching_matrix}}
already uses the Gram adjoint, Schur's lemma, and the norm of the adjoint
on a highest-weight vector. It calibrates that vector using Cartan
multiplication and \code{weylNormSq_le_mul_of_toPoly_eq}.
The readable edition instead uses the extended-tensor argument analyzed below.

The recommendation is to bring the compact proof's representation-theoretic
reduction back to this readable interface. The direct signed-coordinate
or fixed-unit-tensor calibration below gives an alternative that needs
only the existing filling matrix, yields a stronger standalone norm theorem,
and avoids transporting the Cartan norm argument. It is not a claim that
the Schur reduction is newly discovered. Nor is it a measured improvement
over the compact implementation, whose final Cartan calibration is already
very short. In fact the readable source already has the needed Cartan
norm inequality. For the shortest implementation, try the direct port
first; the unit-tensor proof is an optional interface refinement.

\subsection{The general linear-algebra fact}

\begin{theorem}[One-vector calibration of an intertwiner]
\label{thm:onevector}
Let \(W,V\) be finite-dimensional complex Hermitian spaces. Suppose \(B:W\to V\)
intertwines unitary representations, allowing a unitary character twist
on \(V\), and the commutant on \(V\) consists of scalars. Let \(v\in V\)
have norm one. Suppose \(E:W\to\C^J\) is an isometry and \(w\in\C^J\)
has norm one, with
\begin{equation}
 \ip{v}{B\xi}=c\,\ip{w}{E\xi}
 \qquad(\xi\in W).
\end{equation}
Then \(\norm{B\xi}^2\le |c|^2\norm{\xi}^2\) for every \(\xi\in W\).
No multiplicity-free decomposition of \(W\) is assumed.
\end{theorem}

\begin{proof}
Let \(D=B^\dagger\). Equivariance and unitarity imply that \(BD\)
commutes with the target representation. By the commutant hypothesis,
\(BD=\lambda I_V\). Since \(BD\) is positive semidefinite,
\(\lambda\) is real and nonnegative.

Equation (9.1) and Cauchy--Schwarz give
\[
 |\ip{Dv}{\xi}|=|\ip{v}{B\xi}|
 \le |c|\norm{\xi}.
\]
Take \(\xi=Dv\); treating the zero case separately yields
\(\norm{Dv}\le|c|\). Hence
\[
 \lambda=\ip{v}{BDv}=\norm{Dv}^2\le|c|^2.
\]
For arbitrary \(\eta\in V\),
\(\norm{D\eta}^2=\lambda\norm{\eta}^2\).
For \(z=B\xi\), Cauchy--Schwarz now gives
\[
 \norm z^2=|\ip{Dz}{\xi}|
       \le\sqrt{\lambda}\norm z\,\norm\xi.
\]
If \(z=0\) the conclusion is immediate; otherwise divide and square.
\end{proof}

Equivariance makes all target directions have the same adjoint norm.
Only one target direction needs a coefficient calculation.

\subsection{The source already supplies the hypotheses}

Take \(B\) to be the integral branching matrix in
\src{Lib/GLn/Branching/NormBound}{\code{weylNormSq_mulVec_le_of_equivariant}}.
The two Gram matrices \(G_W,G_V\) define the source and target Hermitian
forms. The existing filling matrix satisfies
\begin{equation}
 G_W=E^*E,
 \qquad E=\operatorname{fillingMatrixOver}(\C,N,\gamma+\kappa').
\end{equation}
Thus \(E\) is an isometry for the source's \code{weylNormSq}.
This is the Gram factorization already used at the end of the current
norm proof.

Choose the target highest-weight evaluation vector
\begin{equation}
 v=\operatorname{weylEval}
   \bigl(n,\kappa_0,\operatorname{permArray}(\operatorname{rowEquivFin})\bigr).
\end{equation}
Its norm is exactly one, by
\src{Lib/GLn/SchurModule/GramMatrix}{\code{weylNormSq_weylEval_permArray}}.
The needed commutant statement is already
\src{Lib/GLn/SchurModule/WeylKernel}{\code{mem_range_scalar_of_commute_schur}},
which applies to tuples of unitary matrices.

The source's hypothesis \code{B_block} is
\begin{equation}
 B\,S_W(g)=\delta(g)S_V(g)B,\qquad
 \delta(g)=\prod_\tau\det(g_\tau)^{\kappa'_\tau(\operatorname{rank}n_\tau)}.
\end{equation}
For unitary \(g\), \(|\delta(g)|=1\). It is exactly the allowed character
twist in the abstract theorem.

One can also stay entirely in the source's matrix formalism. Since \(B\)
and the Gram matrices are real, define
\begin{equation}
 D=G_W^{-1}B^{\mathsf T}G_V.
\end{equation}
This is the Hermitian adjoint. Apply \code{B_block} at \(g^{\mathsf T}\),
transpose, and use
\src{Lib/GLn/SchurModule/WeylKernel}{\code{weylGramOver_mul_schur}} to obtain
\begin{equation}
 S_W(g)D=\delta(g)D S_V(g).
\end{equation}
Combining (9.4) and (9.6) gives
\begin{equation}
 \delta(g)S_V(g)BD=\delta(g)BD S_V(g).
\end{equation}
For unitary \(g\), \(\delta(g)\ne0\), so the existing Schur lemma applies.
This avoids adding an independent representation-theoretic framework.

\subsection{The shortest coefficient proof already has its ingredients}

The readable
\src{Lib/GLn/SchurModule/GramMatrix}{\code{weylNormSq_le_mul_of_toPoly_eq}}
says that the Gram norm is submultiplicative for Cartan multiplication
expressed via \code{toPoly}. The Gram-adjoint identity and \code{B_top} give
\[
 \operatorname{toPoly}_{\gamma+\kappa'}(Dv)
   =c\,K_\gamma(y_{\rm first})K_{\kappa'}(y_{\rm last}),
\]
where \(K\) denotes the corresponding Weyl kernel and the two arrays are
the permutation arrays with the new row first and last.
Each factor is \code{toPoly} of a norm-one evaluation vector.
Apply the existing submultiplicativity lemma and
\code{weylNormSq_smul} to obtain
\(\norm{Dv}^2\le|c|^2\) directly.
This is the compact implementation's calibration, with ingredients
already present in the readable imports.

Combined with the Schur argument above, this is the shortest plausible
source migration. No new Cartan theory is needed. The next subsection
gives a second proof that isolates a reusable single-coefficient
criterion and uses only the filling-matrix isometry.

\subsection{An alternative direct coefficient calibration}

The hypothesis \code{B_top} identifies the polynomial represented by
the row \(v^*G_VB\) with \(c\) times \code{branchVector}.
The lemma \code{branchVector_eq_prod_tableauPoly} writes this polynomial
as the product of the minors specified by \code{branchRows}.
Every column of those rows is an injection.

Sorting each column changes its determinant only by a sign. Consequently
the product is a signed single filling polynomial,
\begin{equation}
 \operatorname{branchVector}=\epsilon\,\operatorname{fillingPoly}(T_0),
 \qquad \epsilon\in\{1,-1\}.
\end{equation}
By
\src{Lib/GLn/SchurModule/TensorModel}{\code{nablaCombination_fillingMatrixOver}}
and injectivity of \code{nablaCombination}, \code{B_top} becomes
\begin{equation}
 \ip{v}{B\xi}=c\,\epsilon\,(E\xi)_{T_0}.
\end{equation}
This is Theorem~\ref{thm:onevector} with a signed coordinate unit vector \(w\).

For implementation, a fixed-tensor proof may be easier than formal sorting.
Define \(w_T\) as the product, over the branch columns, of the corresponding
identity-matrix minors in filling \(T\). Expanding the row substitution
at the identity gives
\begin{equation}
 \operatorname{branchVector}
     =\sum_T w_T\,\operatorname{fillingPoly}(T).
\end{equation}
Cauchy--Binet at the identity, with injective rows, gives
\(\sum_T|w_T|^2=1\).
This proves the same coefficient calibration using one fixed tensor.
It does not require constructing tensors for arbitrary evaluation arrays
or proving their pairwise isometry.

Therefore the exact current conclusion follows:
\begin{equation}
 \boxed{\ 
 \operatorname{weylNormSq}(B\xi)
     \le c^2\operatorname{weylNormSq}(\xi).\ }
\end{equation}
The source has \(c\in\Z\), so \(|c|^2=c^2\).

\subsection{A stronger statement and a concrete deletion target}

The current norm theorem additionally assumes the weight relation
\begin{equation}
 \kappa_{0,\tau}(h)=\kappa'_\tau(h)+\gamma_\tau(h+1)
 \qquad(0<h<\operatorname{rank}n_\tau).
\end{equation}
The new proof of the norm bound does not use (9.12).
The hypotheses \code{B_top} and \code{B_block} already provide the
coefficient and equivariance information. The relation can remain
in callers that construct \(B\), but the norm theorem itself may drop it.
This strengthens the reusable statement while shortening its proof.

The current route constructs \code{extendedTensor}, proves a
Cauchy--Binet identity for arbitrary inputs, divides by determinant
polynomials on a dense open set, proves that the associated evaluation
vectors span, and applies a general projection inequality.

The measured replacement targets in
\src{Lib/GLn/Branching/NormBound}{\code{NormBound.lean}} are 246 physical
lines of helpers starting at \code{extendedTensor}, 41 lines in the final
theorem body, and the four-line \code{leadingMinorAlong} definition.
These are 291 gross candidate lines, not net savings.
Keep \code{branchRows}, \code{branchVector_eq_prod_tableauPoly}, and the
earlier branching construction. New adjoint and fixed-tensor lemmas will
take some space.

A source-wide reference scan found no external uses of the 12 contraction
helpers targeted here. The final norm theorem's external consumer is in
\src{Lib/GLn/Branching/IntegralMap}{\code{IntegralMap.lean}};
keeping the old conclusion makes the change local.

This differs from the existing Gram-adjoint proposal. That proposal changes
the integral and archimedean estimates for a one-lift comparison operator
using an integral Fock Gram identity. Here the operator is the algebraic
branching map, and the replacement uses the target commutant and a single
highest-weight coefficient. It neither changes the Fock integral adjoint
nor assumes its proposed denominator improvement.
